\documentclass[11pt]{article}
\usepackage[T1]{fontenc}
\usepackage{lmodern}
\ifdefined\pdfgentounicode
  \input{glyphtounicode}
  \pdfgentounicode=1
\fi
\usepackage[letterpaper,margin=1.06in]{geometry}
\usepackage{amsmath,amssymb,amsthm}
\usepackage{microtype}
\usepackage[shortlabels]{enumitem}
\usepackage{titling}
\usepackage{xcolor}
\definecolor{linkblue}{HTML}{244D74}
\usepackage[colorlinks,linkcolor=linkblue,citecolor=linkblue,urlcolor=linkblue]{hyperref}
\hypersetup{pdftitle={Hardness of Approximating Realizable Conjunction Decision Lists}}

\newtheorem{theorem}{Theorem}
\newtheorem{lemma}[theorem]{Lemma}
\newtheorem{corollary}[theorem]{Corollary}
\newcommand{\err}{\operatorname{err}}
\newcommand{\OPT}{\operatorname{OPT}}
\newcommand{\QP}{\mathsf{QP}}
\newcommand{\RQP}{\mathsf{RQP}}
\newcommand{\NP}{\mathsf{NP}}
\newcommand{\E}{\mathbb{E}}
\newcommand{\F}{\mathbb{F}}

\title{Hardness of Approximating Realizable\\Conjunction Decision Lists}
\author{}
\date{Draft \enspace\textbullet\enspace November 2026}
\setlength{\droptitle}{-1.8em}
\pretitle{\begin{center}\LARGE\bfseries}
\posttitle{\par\end{center}\vspace{-0.8em}}
\preauthor{}
\postauthor{}
\predate{\begin{center}\normalsize}

\begin{document}
\maketitle

\begin{abstract}
Given an explicit labeled Boolean sample on which a conjunction decision list with
$k$ rules makes no errors, can one efficiently find a slightly longer list that
makes only a fixed fraction of errors? For each fixed $0<\varepsilon<1/2$,
we rule out deterministic quasipolynomial-time algorithms that guarantee error at
most $\varepsilon$ and length at most
$k\exp(c_\varepsilon(\ln(N+2))^{1/6})$, for some $c_\varepsilon>0$, unless
$\NP\subseteq\QP$. The sample has $N$ equally weighted indexed examples and
dimension $d\le\log_2 N$. The result covers conjunctions of arbitrary width and
arbitrary first-match order. The reduction uses adjacent slices of a sum of
bounded integers. A coupling changes one coordinate group. On that group a local
sunflower argument retains the first matching rules while leaving the opposite
endpoint's address bits unexposed. We also obtain the randomized bound under
$\NP\nsubseteq\RQP$, where $\RQP$ denotes one-sided randomized quasipolynomial time.
\end{abstract}

\section{Introduction}

A conjunction decision list \cite{Rivest} is an ordered sequence of rules. Each rule tests a
conjunction of signed Boolean literals and predicts a bit. The first matching
rule determines the prediction, with one default bit used when no rule matches.
Rule order can make a conjunction useful even if it matches examples of both
labels. This makes approximation by a short decision list different from covering
a sample with rules that are individually correct on every match.

Our input is an explicit indexed sample $S=((x_j,y_j))_{j=1}^N$, with
$x_j\in\{0,1\}^d$ and $y_j\in\{0,1\}$. The sample is \emph{consistent} if equal
feature vectors have equal labels. Repeated vectors are allowed and each index
contributes equally to empirical error
\[
  \err_S(L)=\frac{1}{N}\sum_{j=1}^{N}\mathbf 1\{L(x_j)\ne y_j\}.
\]
The length $|L|$ counts paid rules. The final default is free.
Let $\OPT_k(S)=\min_{|L|\le k}\err_S(L)$. The algorithm receives $S$ and $k$.
It may test any conjunction of the input literals. We use $\QP$ for
deterministic time $2^{(\log(v+2))^{O(1)}}$ on inputs of length $v$. The class
$\RQP$ has the same worst-case time bound with one-sided error. Its algorithms
always reject NO inputs and accept YES inputs with probability at least $2/3$.
We write $[a]=\{1,\ldots,a\}$ for a positive integer $a$.

\begin{theorem}\label{thm:main}
For every fixed $0<\varepsilon<1/2$ there is a constant $c_\varepsilon>0$
with the following property. Suppose an algorithm, on every consistent indexed
sample with $d\le\log_2N$ and $\OPT_k(S)=0$, returns a conjunction decision list
$L$ satisfying
\[
  \err_S(L)\le\varepsilon,\qquad
  |L|\le k\exp\!\left(c_\varepsilon(\ln(N+2))^{1/6}\right).
\]
A deterministic quasipolynomial-time algorithm with this guarantee implies
$\NP\subseteq\QP$. A randomized quasipolynomial-time algorithm that meets both
guarantees with probability at least $2/3$ on every promised input implies
$\NP\subseteq\RQP$.
\end{theorem}

For every fixed $A,C>0$ the permitted length exceeds $Ak(\log(N+2))^C$
for sufficiently large $N$. At error $1/2$, an empirical majority default
suffices. At $d\le\log_2N$, all $3^d$ conjunctions can be enumerated in
polynomial time. The theorem concerns the selection and ordering of rules
into a short list. It covers arbitrary indexed supports, where
multiplicity matters. The reduction uses this freedom to realize its
distribution by explicit enumeration.

\subsection{Related work}

\paragraph{Minimum consistent decision lists.}
Hancock, Jiang, Li, and Tromp \cite[Section 4]{HJLT} prove that approximating
the minimum length of a perfectly consistent conjunction decision list is
NP-hard within a factor $d^\gamma$ for some $\gamma>0$. Their reduction has
one coordinate per vertex. For each vertex there is a positive example with
just that coordinate zero. For each edge there is a negative example with
just its endpoints' coordinates zero. The length
gap relies on exact agreement. In fact, assign each vertex an independent
uniform color in $[a]$ and give each color class the conjunction of all
positive literals outside it. The resulting DNF accepts every vertex example.
It accepts a negative edge example exactly when both endpoints receive the
same color. Even with arbitrary row weights there is an $a$-term DNF with
error at most $1/a$. Nock and Jappy \cite{NJ} study exact consistency under
several size measures, and Ordyniak et al.\ \cite{OPRS} give parameterized
algorithms and hardness results for exact fitting, including decision lists.

\paragraph{Realizable, fixed-error learning.}
Khot and Saket \cite{KS} show hardness of weakly learning a two-term DNF
using any fixed number of DNF terms. Ghoshal and Saket \cite[Theorem~1.1]{GS}
extend the soundness to arbitrary Boolean combinations of any fixed number of
halfspaces. Since each conjunction is a halfspace, their theorem also covers
conjunction decision lists of any fixed length. Their exact YES witness has
two terms. When $d=O(\log N)$, the
at most $3^{2d}$ pairs of conjunctions can be checked in polynomial time.
Our witness length grows with the instance and the NO length bound grows
faster than every fixed power of $\log N$.

\paragraph{Decision trees and bicriteria algorithms.}
Koch, Strassle, and Tan \cite{KST23,KST24} prove fixed-error and size-approximation
hardness for proper decision-tree learning. A tree can be expressed as a list
with one conjunction per leaf, so a lower bound on tree size alone leaves the
list side of the gap to be proved. Frost et al.\ \cite[Theorem~3]{FLMM}
give length, coverage, and accuracy guarantees for lists that may abstain
and hand off to an external classifier. Their running time depends on an
expanded family of tests that can contain a
predicate of the form $r_j\wedge\bigwedge_{i<j}\neg r_i$. Such a predicate may
require several conjunctions of input literals. Theorem~\ref{thm:main} treats
total prediction with a single bit default and measures the number of
conjunctions used in the output.

\paragraph{Decision-list compression.}
Lovett, Wu, and Zhang \cite{LWZ} give size bounds for approximating
bounded-width decision lists under the uniform cube. In our proof we need
a subsequence that preserves the \emph{identity} of the first match after
a local restriction. Its guarantee uses independence only among the variables
of retained tests. This lets us select the tests before knowing the
projection that identifies raw field indices. To our knowledge,
Theorem~\ref{thm:main} is the first hardness result for unrestricted
conjunction decision lists to combine a zero-error witness, fixed additive
error, and a superpolylogarithmic length gap at logarithmic dimension.

\subsection{Proof overview}

Each sample row contains multiple left--right pairs. Only one pair carries the
label; the others are independently inactivated. The active pair has $t$ groups.
In each group a projection-game edge specifies which two rank fields share a
random jitter. The two modes set the difference of the sums of baseline ranks
to zero or one. A satisfying labeling picks one field at each endpoint in each
group, so the jitters cancel and a rank-ordered list classifies exactly.

For soundness we compare the two modes by changing one baseline in a random
group. Independent masks bound how many fields and opposite address bits a
matching rule can test. We freeze the other groups and filter the two
lists anchored at the active addresses to bounded subsequences. The filters
leave the opposite address in the changed group symbolic, so the projection
game still controls collisions among their selected field indices. We couple
the ranks named by these subsequences and charge possible changes only to the
tests of the actual first matching rules. A separate argument using the
inactive product distribution controls rules without a unique active anchor.

\section{The projection game and addresses}\label{sec:source}

A projection game consists of a bipartite multigraph and one map from the
left alphabet to the right alphabet per edge occurrence. Its value is the
maximum fraction of occurrences on which the endpoint labels satisfy the
projection. We use a source with constant alphabets, perfect completeness,
arbitrarily small fixed NO value $\delta$, and the same positive fiber size
$b$ for every right label on every edge. Here and below all edge sampling
is uniform by occurrence.

Such a source follows from the classical gap for
Max-E3-SAT \cite[Theorem~6.5]{Hastad} and Raz's parallel repetition theorem \cite{Raz}.
Use the gap for clauses on exactly three variables, with satisfiable YES
formulas and NO value bounded away from one by an absolute constant.
In the clause--variable game the left question is
a clause and the right question is one of its three incident variables.
On the left, use the seven satisfying assignments to the clause's three
\emph{literal truth values}, plus a second copy of each of the three
assignments with exactly one true literal. This is a common alphabet of
ten labels. For each position, precisely five of the labels give each
truth value. The projection converts the assigned literal truth value
to the underlying variable value, accounting for negation. This is an
onto five-to-one map, as in Wenner's regularization \cite[proof of
Theorem~1.17]{Wenner}. A global satisfying assignment wins every edge.
For any labeling in NO, each clause violated by the right-hand assignment
forces a disagreement on at least one of its three edges. Thus the
game has value at most $1-\gamma_0$ for some absolute $\gamma_0>0$.
Repeat $r$ independent edges in parallel, requiring all $r$ projections
to pass. For a suitable constant $r$ the repeated game has value at
most $\delta$ and fiber size $b=5^r$, with a polynomial number of
occurrences in the SAT input length. Equivalently, the perfect-completeness
four-to-one theorem of Fei, Minzer, and Wang \cite[Theorem~1.6]{FMW}
supplies a source with $b=4$.

\subsection{Smoothness and semantic addresses}

Discard isolated vertices and take enough disjoint copies that the
source has at least $v$ occurrences, where $v$ is the input length
of the SAT instance. We construct a new game for one group.
Draw $h$ independent source edges $(u_j,w_j)$, $1\le j\le h$, and an
independent uniform position $i\in[h]$. Give the left player
$(u_1,\ldots,u_h)$ and the right player $(i,w_i,(u_j)_{j\ne i})$.
A left label assigns one original left label at each position. The right
label uses an original right label at position $i$ and original left
labels at all other positions, in their induced order. Apply the original
projection at $i$ and copy all other answers. The new map is onto
and still has fiber size $b$. A satisfying labeling extends. In NO,
conditioning on $i$ and on the other left vertices yields a pair of
original-game strategies for coordinate $i$. The new value is at
most $\delta$. Also, for a fixed left vertex and two distinct left
labels, a collision is possible only when they differ at exactly
one position and that position is $i$. Its conditional probability
is at most $\gamma=1/h$ because $i$ is independent of the left vertex.

We can increase the right alphabet size $R$ arbitrarily by
appending a common copy index to the two alphabets and preserving it
under every map. This preserves the fiber size and completeness.
Forgetting indices sends each successful edge of a labeling to a
successful edge of a labeling of the preceding game, so NO value
does not increase. Labels with different copy indices never collide.
Let $q=bR$. Both alphabets are fixed independently of $v$. The group
game is explicit, with $M=v^{O(1)}$ occurrences and $M\ge v$.

Give each vertex of occurrence degree $d_u$ a set of $d_u$ copies.
On the sampled occurrence, draw the two degree copies independently and
uniformly conditional on their vertex. Also choose independent uniform
permutations $P$ of $[q]$ and $T$ of $[R]$ at the left and right endpoints.
Add independent uniform dummy indices with $R!$ choices on the left and
$q!$ on the right. The raw projection on this occurrence is
\[
  \rho=T^{-1}\pi P ,
\]
where $\pi$ is the group game's projection. The original vertex, degree
copy, permutation, and dummy index together form a \emph{semantic address}.
Each side is uniform on its set of semantic addresses, since an endpoint's
degree factor in its marginal cancels the uniform copy choice. There are
$M q!R!$ such addresses on each side. Conditional on its original vertex,
the randomness of an endpoint's own semantic address is independent
of the underlying edge.

\subsection{Balanced address encoding}\label{sec:codes}

Set $t=\lceil\log_2(v+2)\rceil$ and $n=\lceil C_n t\rceil$,
where $C_n$ will be a sufficiently large constant. Choose nested linear
spaces $D\subset F\subset\F_2^n$ of dimensions $\lfloor n/4\rfloor$
and $\lfloor n/2\rfloor$. For an absolute constant $\alpha>0$ we can
require both $F$ and $D^\perp$ to have minimum distance greater than
$2\alpha n$. To see this, choose the nested pair uniformly at random.
The two spaces have uniform marginals at their dimensions. A fixed
nonzero vector belongs to $F$ with probability $O(2^{-n/2})$ and
to $D^\perp$ with probability $O(2^{-n/4})$. For small enough $\alpha$,
the number of vectors of weight at most $2\alpha n$ is
$2^{(h_2(2\alpha)+o(1))n}$, where $h_2$ denotes binary entropy and
$h_2(2\alpha)<1/4$. A union bound proves existence for all sufficiently
large $n$. Enumerating generators and checking distances takes
$2^{O(n^2)}$ deterministic time.

Encode each semantic address by a distinct coset of $D$ in $F$,
separately for each side. There are enough cosets for a sufficiently
large $C_n$ because $M$ is polynomial in $v$. Conditional on an edge
and its two semantic addresses, independently sample a uniform word in
each of the assigned cosets. Any set of at most $\alpha n$ coordinates
of either sampled word is jointly uniform on its Boolean cube.
Indeed, failure of surjectivity of the projection of $D$ onto those
coordinates would give a nonzero annihilator in $D^\perp$ supported there.

Publish each bit $x$ as the pair $(x,1-x)$. Concatenate the $t$ group
words to make an endpoint's \emph{full address}. A full address is valid
if every group word belongs to the assigned coset of some semantic
address of that group and side. On either side the full address is uniform
on its set of $z$ valid addresses, with $\log_2 z=\Theta(tn)$.
Each has exactly $tn$ zero bits. Two different valid full addresses
are at relative Hamming distance at least $\alpha/t$.

\section{The sample and its exact witness}\label{sec:construction}

Let \(\zeta>0\) be a small constant to be fixed from \(\varepsilon\).
Use
\begin{equation}\label{eq:m}
  m=\left\lceil C_m t^2(1+\log_2 z)\right\rceil
\end{equation}
pairs in a sample row, with $C_m$ a sufficiently large constant.
Each pair contains two endpoints. An endpoint stores its full address,
one activation bit, and raw rank fields for all $t$ groups.

Fix a large power of two $H$, let $Q=H-5$, and put
$B_H=\lceil(\log_2 H)^2\rceil$. A pair \emph{active in mode $y$}, for
$y\in\{0,1\}$, has the following law.

\begin{enumerate}
\item Independently in each group draw a game occurrence and the two
addresses described in Section~\ref{sec:source}. Set both activation bits
to one.
\item Draw $B_{Li},B_{Ri}\in\{0,\ldots,5\}$ for $i\in[t]$ uniformly
conditional on
\begin{equation}\label{eq:baseline}
  \sum_{i=1}^t(B_{Li}-B_{Ri})=y.
\end{equation}
These are the baselines. Independently of them, in each group $i$
associate to each raw right index $a$ a jitter $J_{ia}$, independently
over all groups and right indices. Its probability at $j\in\{0,\ldots,Q-1\}$
is proportional to
\[
    (1+1/B_H)^{\min\{j,Q-1-j,\lfloor Q/10\rfloor\}}.
\]
The two endpoint atoms have the same mass $p$, which satisfies
\begin{equation}\label{eq:taper}
  p\le \exp\big(-\Omega(H/(\ln H)^2)\big).
\end{equation}
This follows by comparing the endpoint weight to a plateau weight.
Every atom is positive, and consecutive atom ratios belong to
$[(1+1/B_H)^{-1},1+1/B_H]$.
\item For the right field $a$ in group $i$, the rank is $B_{Ri}+J_{ia}$.
For the left field $f$ it is $B_{Li}+J_{i,\rho_i(f)}$. Thus all ranks are
in $\{0,\ldots,H-1\}$. For each raw field, independently sample
$U$ uniformly in $\{0,\ldots,H-1\}$ and publish the two $\log_2 H$-bit
words for $U$ and $V=(U+\mathrm{rank})\bmod H$. We call these two words
a packet. Conditional on the rank, each of its words is individually
uniform. The rank is decoded as $(V-U)\bmod H$.
\end{enumerate}

An \emph{inactive pair} is generated from a pair active in a fair random
mode. At each endpoint independently, with probability $1/2$ clear the
activation bit, and otherwise select a uniform zero address bit and
raise it to one. The raise increases the weight of the address and
therefore makes it invalid. The rank packets stay fixed. A sample row
serves one uniform pair index, generates that pair active in a fair
mode $y$, and independently makes all other pairs inactive. Its label
is $y$. We first analyze this ideal law.

\begin{lemma}\label{lem:yes}
The labeled law is consistent. In YES it admits a zero-error list with
\[
  k=2mzH^{2t}
\]
paid rules and one default.
\end{lemma}
\begin{proof}
The served pair is identified by two valid full addresses and
activation bits equal to one. Decode its packets. The difference
between the mean left rank and mean right rank in group $i$ is
$B_{Li}-B_{Ri}$ because every jitter appears exactly $b$ times
on the left and once on the right. Summing gives $y$ by
\eqref{eq:baseline}. Thus the label is a function of the feature vector.

In YES, choose a satisfying labeling of the group game. On a given
address the selected raw index in each group is the inverse-permuted
label of its semantic vertex. At each pair, side, and valid full
address, place one rule for each of the $H^{2t}$ assignments to the
selected $t$ packets. The rule tests those packets, the exact full
address, and activation one. Both endpoints of the served pair
supply exactly one match, and the inactive endpoints supply none.
Sort the entire list by increasing sum of the decoded ranks of a
rule, breaking ties in favor of the left side. Left rules output zero
and right rules output one. The projections make the selected
jitter sums agree. The sums of selected ranks therefore tie in
mode zero and the right sum is smaller by one in mode one.
\end{proof}

\section{Local tools for the soundness proof}\label{sec:tools}

\begin{lemma}[Coupling adjacent baseline slices]\label{lem:adjacent}
There is a joint draw of two baseline vectors whose marginals are
within $o(1)$ in total variation of the two uniform slices in
\eqref{eq:baseline}. The vectors differ by a unit change of one
baseline. Conditional on either vector, any specified group is
changed with probability at most $2/t$. The draw is independent
of the group games and addresses.
\end{lemma}
\begin{proof}
Use $5-B_{Ri}$ in place of each right baseline. The slices are now
the vectors of $\{0,\ldots,5\}^{2t}$ with sums $5t$ and $5t+1$.
Pick uniformly an edge of the cube of integer vectors that increments
one coordinate from the first slice to the second. Each endpoint
has at least $t$ eligible coordinates and at most $2t$, so the
conditional choice is uniform over at least $t$ coordinates.

The induced endpoint law weights a vector in its slice by its number
of eligible coordinates. Under independent uniform coordinates this
number is binomial of mean $5t/3$. A Chernoff bound gives deviation
at most $t^{3/4}$ outside probability $\exp(-\Omega(\sqrt t))$.
The central slice has cube mass $\Omega(t^{-1/2})$. The sum has
variance $O(t)$ and a symmetric unimodal mass function, so its central
mass is at least the average mass in an $O(\sqrt t)$ central window
of constant probability. Symmetry and unimodality follow inductively
because convolution with an interval uniform averages over a sliding
window in a symmetric unimodal sequence. Counting the increment edges
using the two degree bounds gives the upper slice at least half the
mass of the lower slice. The degrees therefore also concentrate after
conditioning on either slice. Normalizing the degree weights changes
each slice law by $o(1)$ in total variation, as claimed.
\end{proof}

\begin{lemma}[Retaining first matches]\label{lem:filter}
Consider an ordered list of tests, each a conjunction of nonempty
allowed-set tests on at most $s$ variables. Each variable has at most
$K_0$ possible allowed sets. Give the variables finite reference
marginals whose minimum positive atom is at least $\mu>0$, and assume
the allowed sets intersect their respective reference supports.
For each $\xi>0$ there is a subsequence of size at most
$P_*(s,K_0,\mu,\xi)$ such that the probability that the first matching
index is discarded is at most $\xi$ whenever the variables read by
the subsequence have the product of their reference marginals.
\end{lemma}
\begin{proof}
Delete later copies of identical tests. A sunflower of supports is a
collection with all pairwise intersections equal to the same set, its
\emph{core}. When the allowed sets also agree on the core, call it a
labeled core. Scan the list in order. Discard a test if it contains
the labeled core of $r_*$ earlier retained tests that form a
sunflower. Otherwise keep it. The integer $r_*$ will be chosen below.

For sets of size $a$, a family with more than $a!(r-1)^a$ distinct
sets contains a sunflower of size $r$. To prove this by induction,
take a maximal disjoint subfamily. If its size is less than $r$, its
union of at most $a(r-1)$ elements meets every set. Some element occurs
in more than $(a-1)!(r-1)^{a-1}$ sets. Remove that element, apply
induction on $a$, then add it to the resulting core. The claim at
$a=0$ is vacuous, since at most one such set exists. Apply this count
to each support size with $r=r_*K_0^s+1$ and choose one retained
test per support. Among a sunflower of $r$ supports, more than
$r_*$ tests agree on the labels of the core. The last of those
would have been discarded. With at most $K_0^s$ distinct tests
per support, the total retained count is $O_{s,K_0}(r_*^s)$.

If a discarded index is the first match, one of its witnessing
labeled cores passes and no earlier petal for that core passes.
Use the first completed retained sunflower for each core. It
was completed before every discard it certifies. Conditional on
the core passing, each of its petals passes independently with
probability at least $\mu^s$, including an empty petal which
always passes. Thus the probability of the core passing with
all its petals failing is at most $(1-\mu^s)^{r_*}$.
The number of labeled cores is at most $2^s$ times the retained
count because each core is a subtest of a retained test.
The union bound is at most $\xi$ for sufficiently large $r_*$.
It used only the variables of the retained tests. A discarded
duplicate can never be the first match.
\end{proof}

\begin{lemma}[Crossed permutation queries]\label{lem:crossed}
Let $\pi:[q]\to[R]$ be $b$-to-one, $q=bR$, and let $P,T$
be independent uniform permutations of $[q]$ and $[R]$,
respectively. If $j\in[R]$ is chosen using $P$ and its side's
private randomness, and $i\in[q]$ is chosen using $T$ and
independent private randomness, then
\[
  \Pr\{\pi(P(i))=T(j)\}
  = O\!\left(\sqrt{\frac{b+\log(qR+2)}{R}}\right)
\]
for all sufficiently large $R$ with $b$ fixed.
\end{lemma}
\begin{proof}
Write $p_i$ for the probability of choosing $i$ and fix a threshold
$u>0$. There are at most $1/u$ indices with $p_i\ge u$. Each contributes
at most $1/R$ by fixing the left randomness, dropping the event
that the right side chooses that index, and using uniformity of $T(j)$.
For the other indices, condition on $j$. The right randomness is
still independent of $P$. For each $i,j$ let $a_{ij}$ maximize
$\Pr\{\pi(P(i))=a\mid j\}$ over $a$ (choose an arbitrary value if
$j$ has probability zero). The contribution from the light indices
is at most
\[
  \E_j\sum_{i:p_i<u}p_i\Pr\{\pi(P(i))=a_{ij}\mid j\}
  \;\le\;u\,\E_P\max_j Z_j(P),
\quad Z_j(P)=|\{i:\pi(P(i))=a_{ij}\}|.
\]
Here fixing $j$ and the right randomness determines the right-hand
value $T(j)$, which justifies maximizing over $a$.

Let $\ell=\lceil10\log_2(qR+2)\rceil$. For large $R$ we have
$\ell\le R/2$, so any prescription of projected values on $\ell$
distinct slots has probability at most $b^\ell/(q)_\ell\le(2/R)^\ell$,
where $(q)_\ell$ is a falling factorial. Counting $\ell$-subsets of
matches and taking a union bound over $j$ yields
\[
   \Pr\{\max_j Z_j(P)\ge a\}
   \le R\frac{\binom{q}{\ell}(2/R)^\ell}{\binom{a}{\ell}} .
\]
For $a=\lceil\ell+16eb\rceil\le q$ this is at most $R4^{-\ell}$, using
$\binom{q}{\ell}\le(eq/\ell)^\ell$ and
$\binom{a}{\ell}\ge((a-\ell+1)/\ell)^\ell$.
Since $\max_j Z_j\le q$, this gives $\E\max_j Z_j=O(b+\ell)$,
also when $a>q$. The claimed bound follows by taking
$u=1/\sqrt{R(b+\ell)}$.
\end{proof}

\section{Soundness for arbitrary decision lists}\label{sec:soundness}

Fix a NO instance and a list with $J\le k2^t$ paid rules. A rule may
read literals from any subset of the row. A row's \emph{owner} is the
index of its first matching rule. The final default is
treated as an always-matching rule of index $J+1$.

For a given pair, call a head \emph{charged} if at either endpoint its
address literals allow exactly one valid full address and it tests
activation one there. Assign it one such endpoint and address as its
\emph{anchor}, breaking ties by a fixed rule. This assignment depends
on the syntax of the rule and the set of valid addresses. The same head
may be charged separately for several pairs. An \emph{anchor list}
contains all heads assigned to one anchor, in their original order.
The \emph{charged owner} is the first matching head charged for the
served pair, or a sentinel last in order if none matches.

\subsection{Width and endpoint ranks of a matching rule}\label{sec:width}

The two actual addresses in the served pair are individually uniform on
their valid sets. Since each head contributes to at most one anchor list
for that pair, the expected total size of its two actual anchor lists
is at most $J/z$. For
$D_0=\lceil(J/z+1)e^t\rceil$, the probability that either list exceeds
$D_0$ is $O(e^{-t})$ by Markov's inequality. Call an anchor list of size
at most $D_0$ light. By \eqref{eq:m}, $\log m=o(t)$. Thus
\begin{equation}\label{eq:light-count}
  \log_2(D_0+1)\le 2t\log_2H+O(t).
\end{equation}

For an anchored head, count one variable for each raw packet it tests
in the served pair, and one for each underlying opposite address bit it
tests. Complementary published bits of one underlying bit contribute
one variable. Tests on the own address, the activation bits, and other
pairs do not enter this count. Contradictory address tests can be
discarded. Set $D_*=\lceil10t(1+\log_2 H)\rceil$ and choose $C_n$ in
Section~\ref{sec:codes} so that $D_*<\alpha n$. If the count exceeds
$D_*$, select one literal from each of $D_*$ distinct counted variables.
Conditional on the group occurrences, semantic addresses, full own
address, baselines, and jitters, the selected packet literals are
independent fair tests by the independent masks. The selected opposite
address bits are independent fair tests by the cosets and are
independent of the masks. The head has probability at most $2^{-D_*}$
of matching. We condition on lightness of the own anchor only when
applying this bound. Union over its at most $D_0$ heads, then over the
two sides. By \eqref{eq:light-count}, off $o(1)$ mass every matching
charged head belongs to a light anchor and has count at most $D_*$.
We may therefore filter to this total width by syntax.

We also bound the number of groups where a width-filtered head reads
a field with jitter $0$ or $Q-1$. For any fixed $\lambda>0$, the event
that some such head at a light actual anchor reads an endpoint jitter
in at least $\lceil\lambda t\rceil$ groups costs at most
\[
  2D_0\,2^{D_*}(2p)^{\lceil\lambda t\rceil}=o(1)
\]
for sufficiently large fixed $H$. For each head choose one such field
in each specified group, using at most $2^{D_*}$ subsets of fields.
The respective jitters are independent across distinct groups
conditional on the occurrences and addresses, with total endpoint
mass $2p$ in each group. The bound follows from \eqref{eq:taper}
and \eqref{eq:light-count}.

Use the coupling of Lemma~\ref{lem:adjacent} and let $g$ be its changed
group. We will draw the remaining data with its prescribed conditional
law in each mode, even conditional on the baseline pair and $g$.
It follows that the conditional probability bound $2/t$ for $g$ remains
valid given either entire latent marginal row. Put
\[
  s=\left\lceil\frac{24(1+\log_2 H)}{\zeta}\right\rceil,
  \qquad \lambda=\zeta/4 .
\]
A width-$D_*$ head can exceed local width $s$ in at most $D_*/s$
groups. Off the endpoint-jitter exception above, it tests endpoint
jitters in fewer than $\lceil\lambda t\rceil$ groups. By the $2/t$
bound, the charged owner in either row, if it is paid, has the
following properties except on at most $2\zeta+o(1)$ mass: it is
at a light anchor, satisfies the total width bound, uses at most $s$
counted variables in group $g$, and reads no endpoint jitter there.
The same width and endpoint estimates apply to the biased marginal
baselines, at an extra $o(1)$ cost.

\subsection{Filtering with the opposite address unexposed}\label{sec:selection}

Share the served index, all inactive pairs, and the occurrences and
addresses of the active pair across the two modes. Share packets in
all groups except $g$, where the baseline can change. Fix the outside
data, including the other groups' addresses and packets, and fix the
baseline pair and $g$. The changed group's game and address randomness
remain fresh. In each of the two anchor lists, drop a heavy own anchor,
apply the total width bound and the local bound $s$, and evaluate all
tests outside the group. The own full address and both activation bits
are fixed for that list. Leave the opposite endpoint's changed
address segment as variables. In particular the symbolic list for a
given side can be formed from its own full address and the fixed
outside data before observing the opposite segment or the local
packets.

For each of the 36 ordered choices of the two local baselines, apply
Lemma~\ref{lem:filter} with $\xi=\zeta/2$. The reference variables are
independent raw packets with the relevant single-packet law and
independent fair opposite address bits. The minimum positive atom is
at least $\min(p/H,1/2)$ and the number of possible unary tests is
bounded in terms of $H$. A unary test with zero reference mass can
never pass under the actual one-variable marginal either, so such a
head can be removed. Take the union of the 36 bounded subsequences
for each side. Each union mentions at most $K=36sP_*$ raw slots and
address bits. For the fixed outside data, the union of a given anchor
depends on that anchor and the local baseline choices, with no
dependence on the opposite local address. The bound $K$ is fixed
before choosing the game parameters.

We seek two properties of the selected slots. Within each anchor's
union, distinct left slots should project to distinct classes. In the
union of both sides' lists, each class should be tested on at most one
side. Distinct right slots already index distinct classes. We bound
the probability of violating these properties conditional on the
outside data, which leave the changed group's source fresh. A slot
ordinal in an anchor's union is a function of its own address.
For comparisons between the two lists we may fill missing ordinals
arbitrarily. The following estimates apply to each relevant pair.

\begin{enumerate}[(i)]
\item Two distinct left slots selected by the left anchor collide with
probability at most $\gamma$. Fix that anchor, so its two left labels
after permutation are fixed and the edge is still random conditional
on the underlying left vertex. For two distinct left slots selected
by the right anchor, condition on the edge and right anchor. The
fresh left permutation gives probability $(b-1)/(q-1)$.
\item A left slot chosen by the left anchor and a right slot chosen by
the right anchor collide with probability at most $\delta$. Undo
the permutations. The two choices give strategies for the group game
using their own vertices and independent private randomness. The NO
value also bounds such strategies by averaging over their private
randomness.
\item For opposite-side slots chosen by the same anchor, condition on
the edge and that anchor. The opposite permutation is uniform and
independent of the choices. The regular fibers give probability $1/R$.
\item For a left slot chosen by the right anchor and a right slot
chosen by the left anchor, condition on the underlying edge.
The two independent permutations and private randomness satisfy
Lemma~\ref{lem:crossed}.
\end{enumerate}

Choose $\delta,\gamma$ sufficiently small and then increase $R$
with the copy index, keeping $b$ fixed. The $O(K^2)$ comparisons
then fail with total probability at most $\zeta$.
In (i) the requirement of two distinct existing slots is measurable
at the selecting anchor. The bounds still apply when only some anchors
contribute the two ordinals. Same-side overlap between different
lists is allowed.

To bound the filtering loss of \emph{one} anchor, condition on the
edge, both semantic addresses, the full own address, the baseline pair
and the outside data. Injectivity of this anchor's retained slots
across its different raw fields is determined by this information:
it requires distinct classes among the left slots and between the
left and right slots in that union. Under this injectivity, its
retained packets have independent product marginals. The opposite
address coordinates used in the retained tests are jointly fair
since $K<\alpha n$ for sufficiently large $v$, and they are independent
of the packets. Apply Lemma~\ref{lem:filter} under this conditioning.
The probability of a first match discarded by this anchor's filter,
intersected with its own injectivity, is at most $\zeta/2$ for either
mode. This calculation is made separately for each anchor.
Conditioning on the other anchor's union could reveal the very
opposite address bits being kept uniform.

If the full charged owner is eligible for its local anchor list,
it is that local list's first match. On the event of no collisions,
the probability of losing such an owner in the union is at most
$\zeta$ per mode by summing the two filter loss bounds. This holds
even if the full list has many matching heads.

\subsection{Coupling the actual charged owners}\label{sec:coupling}

To sample the changed group, now fix both full addresses, the edge,
the outside data and both baselines. The retained unions and the
collision event are determined. If a collision occurs, use any
coupling with the prescribed marginals. Otherwise each selected
jitter class has a unique used side. For every such class independently
maximally couple its two ranks on that side, and subtract the
corresponding baselines to obtain the two jitters. For unselected
classes share their jitters, and use a common independent uniform
$U$ in both rows at every raw field, independently across raw
fields. This gives the specified conditional product law for each
marginal. In particular it verifies the latent row claim used in
Section~\ref{sec:width}.

On an unchanged side the rank coupling preserves the rank. On a
side whose baseline changes, a rank with nonendpoint jitter lies
in both rank supports. At that rank the ratio of the two masses
is between $(1+1/B_H)^{-1}$ and $1+1/B_H$. In a maximal coupling
the common mass at each rank is the smaller of its two masses.
Thus, conditional on that rank in either row, the chance it
changes is at most $1/B_H$. This remains true conditional on
that row's full latent data and the shared frame, since the local
couplings are independent across classes and were chosen before
drawing jitters and masks. If the rank is preserved, both packet
words are preserved.

For a retained actual charged owner of local width at most $s$, with
no endpoint jitter in this group, the probability that any of its
tests changes is at most $s/B_H\le\zeta$ for large fixed $H$. This
uses only the tested fields of the owner. If both paid charged
owners still match in the opposite row, they agree by their common
order. The same implication holds when a sentinel is present:
a paid rule preserved on the other row rules out a sentinel there.

Counting the collision event once and the local width and boundary,
filtering, and rank changes for both owners gives a disagreement
probability at most
\[
  \zeta+2(2\zeta+\zeta+\zeta)+o(1)=9\zeta+o(1).
\]
Finally correct each biased baseline marginal to its exact uniform
slice using a maximal coupling applied to its marginal baseline alone.
When its baseline vector changes,
resample its active packets with the correct law. A marginal's
packet law conditional on its old baseline and the common
occurrences and addresses was already the required law. This
correction has total cost $o(1)$ and leaves the common served index,
inactive context, occurrences and addresses fixed. It gives
charged-owner disagreement at most $10\zeta+o(1)$ under the ideal
law.

The parameter choices are noncircular. Fix $\zeta$ and then $H$ so
that the endpoint tail union and $s/B_H\le\zeta$ hold. This fixes
$P_*$ and $K$, so choose $\delta,\gamma$ and the alphabet enlargement
next. Choose $C_n$ large enough for the semantic addresses and
$D_*<\alpha n$. Because $K$ is constant the other code constraint
$K<\alpha n$ holds eventually. The constant $C_m$ affects
\eqref{eq:light-count} only through $\log m=o(t)$.

\section{Restoring the inactive pairs}\label{sec:restoration}

Use the preceding coupling conditional on the served index and all the
other inactive pairs. Let $X^0,X^1$ be the two active rows with exact
marginals. Inactivate the served pair in each row with fresh randomness,
giving $V^0,V^1$. Write $O$ for the first matching index of the full
list, including the default. We first show
\begin{align}
 \Pr\{O(V^0)\ne O(V^1)\}
   &\le \frac{8\ln(J+1)}{m}, \label{eq:restored-agree}\\
 \Pr\{O(V^y)\text{ fails to match }X^y\}
   &\le \frac{2\log_2(J+1)}{m}\quad (y=0,1).
     \label{eq:restored-survive}
\end{align}

\paragraph{Disagreement of restored owners.}
Draw a row from the product of the inactive pair laws and independently
resample one uniform pair. Let $p_j$ be the original owner probability
of rule $j$, and let $a_{ji}$ be the probability its tests in pair $i$
pass. Its owner cell is contained in its match set, so
$p_j\le\prod_i a_{ji}$ and $\sum_i(1-a_{ji})\le-\ln p_j$ whenever
$p_j>0$. Resampling loses the first row's winning match with
probability at most
\[
  \frac{1}{m}\sum_j p_j(-\ln p_j)\le\frac{\ln(J+1)}{m}.
\]
If the owners differ, the earlier owner fails on the other row.
By symmetry of resampling, disagreement has probability at most
$2\ln(J+1)/m$.

Conditional on the chosen index and the rest of the row, let $p_*$ be
the probability of the modal owner under the fair inactive mixture
in that pair. The two independent draws in the resampling experiment
disagree with probability at least $1-p_*$ because their agreement
probability is the sum of squared owner probabilities. The replaced
block of each $V^y$, conditional on the served index and inactive
pairs, has density at most two relative to the same fair inactive
mixture. This follows because it is obtained from just one of the
two fair-mixture modes. In establishing this density bound we
average over the shared served occurrence and addresses.
Therefore each $O(V^y)$ differs from the modal owner with conditional
probability at most $2(1-p_*)$. Averaging and summing over the two
modes proves \eqref{eq:restored-agree}.

\paragraph{Survival of a restored owner on reactivation.}
Generate independent fair-mode active frames in all $m$ pairs and
then independently inactivate all of them. Condition on the active
frames. If a conjunctive rule fails on an active pair, a failed
literal at some endpoint can be repaired by at most one of that
endpoint's one-bit changes. Such a change has probability at most
$1/2$. Let $f_j$ be the number of pairs where rule $j$ fails before
inactivation. By block independence its conditional probability
$p_j$ of becoming the first match is at most $2^{-f_j}$. The
expected number of failed pairs for the restored owner, conditional
on the frames, therefore satisfies
\[
  \sum_j p_j f_j\le \sum_j p_j(-\log_2 p_j)\le\log_2(J+1),
\]
with zero terms omitted. A uniform pair belongs to those failed
pairs with probability at most $\log_2(J+1)/m$. Conditioning
the chosen pair's latent mode to be $y$ multiplies that bound
by at most two. This is exactly the one-mode marginal of $V^y$
with the chosen pair replaced by its active frame in $X^y$,
proving \eqref{eq:restored-survive}.

\paragraph{From charged owners to prediction error.}
Suppose the active owners differ and the earlier one is uncharged.
Conditional on the two active rows, this earlier rule survives
inactivation of its served pair with probability at least
$\alpha^2/(4t^2)$. In fact, at an endpoint where the address tests
allow several valid addresses, let $w'$ be a second valid address
satisfying the tests. By equal weight and relative distance, at
least an $\alpha/t$ fraction of the actual address's zero bits
are ones in $w'$. Each such bit is untested, so raising it
preserves the match. At an endpoint where the tests allow only
one valid address, an uncharged matching rule has no activation-one
test, so clearing the activation bit preserves its match. The
two choices are independent, giving the stated probability.

On a surviving restoration, the restored owner on that side comes
before the other active owner. Thus the two restored owners cannot
both agree and match their respective active rows. Equations
\eqref{eq:restored-agree} and \eqref{eq:restored-survive} show that
the probability of distinct active owners with an uncharged earlier
one is at most
\[
  \frac{4t^2}{\alpha^2}
  \left(\frac{8\ln(J+1)}{m}+\frac{4\log_2(J+1)}{m}\right).
\]
If instead the earlier of two distinct active owners is charged, it
is the charged owner of its own row and fails on the other row.
Hence the charged owners differ. Section~\ref{sec:coupling} bounds
that event by $10\zeta+o(1)$.

Both active predictions being correct implies different owners.
By a union bound the probability both predictions are correct is
at least $1-2\err(L)$, where the error is under the ideal law.
Combining the bounds gives
\begin{equation}\label{eq:no-error}
  \err(L)\ge \frac12-5\zeta-o(1)
                  -O_\alpha\!\left(\frac{t^2\log(J+1)}{m}\right).
\end{equation}
Because $J\le 2mzH^{2t}2^t$ and $\log z=\Theta(t^2)$,
we have $\log(J+1)=O(\log z)$. Choose the constant $C_m$
in \eqref{eq:m} large enough that the last term in
\eqref{eq:no-error} is at most $\zeta$ for all sufficiently
large $v$. Thus every list of length at most $k2^t$ in NO
has ideal error at least $1/2-6\zeta-o(1)$.

\section{The explicit indexed sample and the complexity gap}\label{sec:size}

The ideal law is a finite product of conditional draws. We first
make the degree-copy choices have a uniform integer index of a
size independent of the drawn edge. For an occurrence count $M$
and requested degree $d_u$, use an independent uniform integer
in $\{0,\ldots,M^3-1\}$ reduced modulo $d_u$. Every copy occurs,
and its probability differs from uniform by at most $1/M^3$.
Since $d_u\le M$, the variation cost per draw is $O(M^{-2})$.
The total cost over all pairs and groups is $O(mt/M^2)=o(1)$
because $m,t$ are polylogarithmic in $v$ and $M\ge v$.

All other choices admit uniform indices directly. The positive
rational jitter weights can be scaled to fixed integers.
For the baseline slices, use the product of their cardinalities
as a common denominator, which takes $O(t)$ bits since each
slice has at most $6^{2t}$ vectors. The fair mode and the
inactivation step also have fixed denominator sizes for a given
instance, with exactly $tn$ possible zero bits at each endpoint.
Unused star choices at the served pair can still be enumerated.
The edge occurrences, permutations, dummy indices, address cosets,
and packet masks are explicit finite sets. Enumerate all indices
of the resulting uniform product. Repeat a row whenever it has
several indices. This preserves the ideal support, so consistency
and the YES witness are exact by Lemma~\ref{lem:yes}. Its variation
distance from the ideal law is $o(1)$, so in NO it gives the same
error gap up to $o(1)$.

An enumeration index for one pair takes $O_\varepsilon(t^2)$ bits.
Each of its $t$ groups uses $O_\varepsilon(\log(v+2)+n)$ bits for
the occurrence, address and copy choices and a constant number
of bits for the rank law. The slice and star indices fit the same
bound. From \eqref{eq:m} and $\log z=\Theta(t^2)$, we have
$m=O_\varepsilon(t^4)$. The dimension before padding is at most
$O_\varepsilon(mt(n+1))=O_\varepsilon(t^6)$.
By uniformly duplicating the indices if needed, ensure $N\ge2^d$.
This keeps the same empirical law and gives
\begin{equation}\label{eq:clock}
  \ln(N+2)=O_\varepsilon(t^6),\qquad
  t=\lceil\log_2(v+2)\rceil .
\end{equation}
Constructing the nested codes by search, enumerating the indices,
and padding the sample all take deterministic quasipolynomial time.

\begin{proof}[Proof of Theorem~\ref{thm:main}]
Fix $\zeta<(1/2-\varepsilon)/20$ and a rational
$\tau\in(\varepsilon,1/2-6\zeta)$.
For all sufficiently large $v$, the deterministic construction
maps SAT to the following promise gap for a consistent indexed
sample with $d\le\log_2N$.
\begin{itemize}
  \item \textbf{YES.} A list with $k$ paid rules has zero error.
  \item \textbf{NO.} Every list with at most $k2^t$ paid rules
        has error greater than $\tau$.
\end{itemize}
Choose $c_\varepsilon>0$ sufficiently small that
$c_\varepsilon(\ln(N+2))^{1/6}\le(\ln2)t/2$ by \eqref{eq:clock}.
The learner's permitted length is then at most $k2^t$ for
large inputs, even after any fixed additional prefactor.

Run the supposed algorithm with its quasipolynomial clock enforced
also outside the realizable promise. Accept only when the output
has valid syntax, at most $k2^t$ paid rules, and empirical error at
most $\tau$. These conditions are checked by evaluating the first
match on every indexed example. On YES the guarantee supplies an
accepted output. On NO the gap rejects every output. The finitely
many smaller SAT inputs can be decided by exhaustive search.
This gives deterministic quasipolynomial time in $v$ for a
deterministic learner and the same worst-case running time with
one-sided error for a randomized learner. In the latter case
the guarantee supplies an accepted output with probability at
least $2/3$ on YES, and every coin outcome is rejected on NO.
Composing with polynomial-time reductions to SAT proves the
claimed implications.
\end{proof}

\begin{corollary}\label{cor:coins}
A randomized quasipolynomial-time learner as in
Theorem~\ref{thm:main} that uses at most
$(\log(N+2))^{O(1)}$ random bits in the worst case implies
$\NP\subseteq\QP$.
\end{corollary}
\begin{proof}
Enumerate all random tapes up to the stated bound and use the
same verifier. A YES input has a passing tape, and every tape
fails on NO. By \eqref{eq:clock} the tape enumeration and
verification take deterministic quasipolynomial time in $v$.
\end{proof}

\begin{thebibliography}{99}

\bibitem{Rivest}
Ronald L. Rivest.
\newblock Learning decision lists.
\newblock \emph{Machine Learning}, 2(3):229--246, 1987.
\newblock \href{https://people.csail.mit.edu/rivest/pubs/Riv87b.pdf}{Author copy}.

\bibitem{HJLT}
Thomas R. Hancock, Tao Jiang, Ming Li, and John Tromp.
\newblock Lower bounds on learning decision lists and trees.
\newblock \emph{Information and Computation}, 126(2):114--122, 1996.
\newblock \href{https://doi.org/10.1006/inco.1996.0040}{DOI}.

\bibitem{NJ}
Richard Nock and Pascal Jappy.
\newblock Decision tree based induction of decision lists.
\newblock \emph{Intelligent Data Analysis}, 3(3):227--240, 1999.
\newblock \href{https://doi.org/10.1016/S1088-467X(99)00020-7}{DOI}.

\bibitem{OPRS}
Sebastian Ordyniak, Giacomo Paesani, Mateusz Rychlicki, and Stefan Szeider.
\newblock A general theoretical framework for learning smallest interpretable models.
\newblock \emph{Artificial Intelligence}, 350:104441, 2026.
\newblock \href{https://doi.org/10.1016/j.artint.2025.104441}{DOI}.

\bibitem{KS}
Subhash Khot and Rishi Saket.
\newblock Hardness of minimizing and learning DNF expressions.
\newblock In \emph{FOCS}, pages 231--240, 2008.
\newblock \href{https://doi.org/10.1109/FOCS.2008.37}{DOI}.

\bibitem{GS}
Suprovat Ghoshal and Rishi Saket.
\newblock Hardness of learning DNFs using halfspaces.
\newblock In \emph{STOC}, pages 467--480, 2021.
\newblock \href{https://arxiv.org/abs/1911.06358}{Full version}.

\bibitem{KST23}
Caleb Koch, Carmen Strassle, and Li-Yang Tan.
\newblock Properly learning decision trees with queries is NP-hard.
\newblock In \emph{FOCS}, pages 2383--2407, 2023.
\newblock \href{https://arxiv.org/abs/2307.04093}{Full version}.

\bibitem{KST24}
Caleb Koch, Carmen Strassle, and Li-Yang Tan.
\newblock Superconstant inapproximability of decision tree learning.
\newblock In \emph{COLT}, PMLR 247, pages 2979--3010, 2024.
\newblock \href{https://proceedings.mlr.press/v247/koch24a.html}{Proceedings}.

\bibitem{FLMM}
Nave Frost, Zachary C. Lipton, Yishay Mansour, and Michal Moshkovitz.
\newblock Partially interpretable models with guarantees on coverage and accuracy.
\newblock In \emph{ALT}, PMLR 237, pages 590--613, 2024.
\newblock \href{https://proceedings.mlr.press/v237/frost24a.html}{Proceedings}.

\bibitem{LWZ}
Shachar Lovett, Kewen Wu, and Jiapeng Zhang.
\newblock Decision list compression by mild random restrictions.
\newblock In \emph{STOC}, 2020.
\newblock \href{https://doi.org/10.1145/3357713.3384241}{DOI}.

\bibitem{Hastad}
Johan H{\aa}stad.
\newblock Some optimal inapproximability results.
\newblock \emph{Journal of the ACM}, 48(4):798--859, 2001.
\newblock \href{https://people.kth.se/~johanh/optimalinap.pdf}{Author copy}.

\bibitem{Raz}
Ran Raz.
\newblock A parallel repetition theorem.
\newblock \emph{SIAM Journal on Computing}, 27(3):763--803, 1998.
\newblock \href{https://doi.org/10.1137/S0097539795280895}{DOI}.

\bibitem{Wenner}
Cenny Wenner.
\newblock Circumventing \(d\)-to-1 for approximation resistance of satisfiable
predicates strictly containing parity of width at least four.
\newblock \emph{Theory of Computing}, 9(23):703--757, 2013.
\newblock \href{https://doi.org/10.4086/toc.2013.v009a023}{DOI}.

\bibitem{FMW}
Yumou Fei, Dor Minzer, and Shuo Wang.
\newblock On the hardness of 4-to-1 games with perfect completeness.
\newblock ECCC TR26-179, revision 3, 2026.
\newblock \href{https://eccc.weizmann.ac.il/report/2026/179/}{Report}.

\end{thebibliography}
\end{document}
